% ===== UNIT 32: VeRi-Wild: literature context + reframing (5.10 part 4) =====
% Status: REWRITTEN  (change to REWRITTEN / REVIEWED when done; see CHECKLIST.md)
% Rules: style rewrite only. Keep every \cite, \ref, \label, number, equation and table value as is. Change the wording, not the content.
These numbers should be read against the cross-domain literature, not against models trained on VeRi-Wild. Supervised training on the 30{,}671-identity training set of VeRi-Wild gives 79--89\% mAP on this protocol~\cite{lou2019veriwild,lu2025clipsenet,benmabrouk2026genlimits}. But that is a different regime, because it uses the very labels whose absence transfer measures. The right peers are transfer results. Published UDA methods are transferred from VehicleID, a source with over twenty times more training identities than VeRi-776, and they also adapt on unlabelled VeRi-Wild data. They reach 24.1\% (SPGAN~\cite{deng2018spgan}), 30.0\% (UDAP~\cite{song2020udap}) and 34.7\% (ECN~\cite{zhong2019ecn}) mAP on the same 3{,}000-identity protocol. The strongest reported direct transfer (no target data) reaches 37.7\%, and full viewpoint-aware adaptation reaches 39.7\%~\cite{zheng2020vapc}. Our zero-shot 35.93\% mAP (ViT-B $+$AQE) comes from the much smaller VeRi-776 source, and touches no VeRi-Wild data. It sits inside this band, above several methods that do adapt on the target. It stays inside the band at the 5{,}000 and 10{,}000 protocols (30.5 and 23.8 vs.\ 30.6 and 24.7 for the adapted ECN). This suggests that the transferring stream is competitive with cross-domain practice at every scale, and is not an outlier.

This cross-dataset probe puts our contributions in a more honest light. The reproducible ViT-B retrieval pipeline generalises in part. The small VeRi-776 advantage of the fusion is domain-specific, and it should not be read as a transferable gain. We argue that reporting a benchmark where the fusion is worse than the stronger stream is more useful to the community than a curated in-domain number.
